%!TEX root = ../main.tex
% =============================================================================
%  摘要专用页（格式规范第三条：含标题与关键词，无需翻译成英文，原则上不超过一页；
%              论文自本页开始编页码，页脚中部，阿拉伯数字从 1 起）
%  电子版第一页必须为本页（第十条：承诺书与编号专用页不放入电子版）
%
%  ⚠️ 本页是唯一有硬性「一页」限制的页面。每改一次都要重新确认：
%     编译后本页仍为 1 页（用 \draftfalse 定稿模式再确认一次）。
%
%  ⚠️ 摘要六条铁律（写作时逐条核对）
%     1 每个数字都必须真算出来，未算出写 {待计算}，严禁编造
%     2 必须有量化结果，不能只写"建立了模型并求解"
%     3 必须有一句"为什么这样建模"（不是模型名罗列）
%     4 必须有一句验证
%     5 必须有一句适用边界
%     6 篇幅控制在一页内
%
%  ⚠️ 写作顺序：最后写摘要，但先写摘要骨架——建模一开始就把下面的占位符抄下来，
%     每算出一个数就填进去。摘要数字 = 正文数字 = 表格数字 = 结果文件数字，零差异。
%
%  ⚠️ 篇幅：本页已按"占位符全部为【待计算】"的最长情形留足余量；
%     填写真实数字后只会更短，不会超页。若增删段落，请重新确认仍为 1 页。
% =============================================================================

% 摘要页行距略紧（1.4 倍），为"占位符全部未填"的最长情形留出余量；
% 格式规范第八条不对行距作统一要求，正文仍为 1.5 倍。
{\linespread{1.4}\selectfont
\begin{center}
  {\zihao{3}\heiti \papertitle\par}
\end{center}

\begin{center}
  {\zihao{4}\heiti 摘\quad 要}
\end{center}

% ---------------------------------------------------------------- ① 背景与任务
微网在分时电价与光伏出力双重不确定性下运行：光伏盈余集中在电价低谷的午间，
而负载高峰落在电价尖峰的傍晚，储能容量与充放功率均有限。本文按信息结构由弱到强
给出四层输出：确定性最优计划 \(\rightarrow\) 不确定性下的安全边际 \(\rightarrow\)
多阶段滚动调整 \(\rightarrow\) 波动电价下的风险决策。

% ---------------------------------------------------------------------- ② 方法
问题 1 中电价、负载与光伏预测给定，且要求 \(0{:}00\) 与 \(24{:}00\) 储电量相同，
目标与约束均为线性，故建立\textbf{线性规划（LP）}模型，由其\textbf{对偶变量}给出
储电水价值与最优充放电的\textbf{阈值型刻画}。问题 2 的计划依据为日前预报、结算依据
为当日实际值，缺口按 5 倍电价紧急购电，二者时间上分离，故建立
\textbf{两阶段随机规划}模型，以\textbf{报童型临界分位数}刻画安全边际。问题 3 有
\(0{:}00\)、\(6{:}00\)、\(12{:}00\)、\(18{:}00\) 四个预报时刻，故建立含
\textbf{非预期性约束}的\textbf{多阶段滚动优化}模型，并用预报频次对照实验量化
\textbf{信息价值（VoI）}，直接回答"是否需要引入其他时刻预报"。问题 4 中电价本身随机，
故以\textbf{分时骨架加随机残差}两层结构建模电价，用\textbf{分块 Bootstrap} 生成
保相关结构的联合情景，在\textbf{条件风险价值（CVaR）}下重算问题 2、3。

% ---------------------------------------------------------------------- ③ 结果
问题 1 最优全天购电费 \ph~元，相对无储能基线 \ph~元节省 \ph，全天购电量 \ph~kWh，
充、放电量 \ph~与 \ph~kWh，弃光 \ph~kWh。问题 2 在 2025-02-01 至 12-31（334 天）
总费用 \ph~元，其中计划购电费 \ph~元、紧急购电费 \ph~元，紧急购电量 \ph~kWh；
计划覆盖率经验值 \ph~落在报童两界 \ph--\ph~之间。问题 3 引入四个预报时刻后总费用由
\ph~元降至 \ph~元，信息价值 \ph，边际收益在 \(n=\ph\) 后进入平稳带，
故\textbf{不建议}进一步加密预报。问题 4 在波动电价下问题 2、3 的费用分别为
\ph~元、\ph~元，CVaR 前沿表明 \ph。

% ---------------------------------------------------------------------- ④ 验证
全部结果通过\textbf{原始可行性}（功率平衡残差 \(\le 10^{-6}\)~kWh）、
\textbf{对偶最优性}与\textbf{费用独立复算}三重检验，并与无储能、贪心、完美信息三类
基线对照；对效率、运行区间、惩罚倍数、预报误差等 \ph~个参数做灵敏度分析，结论在 \ph~范围内不变。

% ------------------------------------------------------------- ⑤ 结论与建议
在本题参数范围内，储能的核心价值是\textbf{把午间光伏盈余搬运到傍晚高价时段}，
且\textbf{储电量容量}而非充放功率是收益瓶颈；当峰谷价比高于盈亏平衡阈值
\(1/(\eta_c\eta_d)\approx1.235\) 时套利有利。结论适用于价格接受型、不外送、
单一储能的并网微网，允许上网售电或引入需求响应时需重新校准。

\noindent{\heiti 关键词：}\quad 微网；储能调度；线性规划；两阶段随机规划；报童模型；多阶段滚动优化；CVaR；VoI
}
\newpage

% =============================================================================
%  ⚠️ 定稿前必须逐项确认（上面 ⑤ 段中的每一条判断都必须有证据）：
%     · "储电量容量而非充放功率是收益瓶颈"  ← 证据：fig04 边际价值曲线 + 充放功率上限灵敏度
%     · "峰谷价比高于 1.235 时套利有利"      ← 证据：阈值律命题 + 电价灵敏度
%     · "不建议进一步加密预报"                ← 证据：VoI 五组对照 + 加密预报成本侧估计（fig06）
%     · "适用于价格接受型、不外送、单一储能"  ← 与 §2.1 假设 A1、§10.4 推广边界一致
%    任何一条拿不出证据，就把它从摘要里删掉，而不是留着。
% =============================================================================
